\subsection{A hyperk\"{a}hler structure on the coadjoint orbit}
Consider the quaternionic spaces $\Q^m$ and $\Q^n$ viewed as right $\Q$-modules. Then the space $\Hom(\Q^m,\Q^n)$ of $n\times m$ quaternionic matrices is itself a left $\Q$-module, where we denote left multiplication by the quaternions ${\rm i}$, ${\rm j}$, ${\rm k}$
with the operators $I$, $J$, $K$. We can further equip this space with a metric%
\begin{equation}\label{quat_metric}
	g(A,B)=\frac{1}{2}\Tr\bigl(AB^\dagger+BA^\dagger\bigr),
\end{equation}
where $A^\dagger$ is the quaternionic conjugate-transpose. The tuple $(g,I,J,K)$ defines a hyperk\"{a}hler structure on $\Hom(\Q^m,\Q^n)$.
Observe that right multiplication by the compact symplectic group ${\rm Sp}(m)\subset {\rm GL}_m\Q$ preserves the hyperk\"{a}hler structure. The following lemma will facilitate us in switching between different complex structures on this space.
\begin{Lemma}\label{facilitate}
	By writing elements $A$ of $\Hom(\Q^m,\Q^n)$ as
	\begin{equation*}
		A=\begin{cases}
			\displaystyle Q+{\rm j}P^\dagger,\\[1mm]
			\displaystyle\frac{Y^\dagger-X}{\sqrt{2}}-{\rm j}\biggl(\frac{{\rm i}\overline{X}+{\rm i}Y^{\mathsf{T}}}{\sqrt{2}}\biggr),\vspace{1mm}\\
			\displaystyle\frac{U+{\rm i}V^\dagger}{\sqrt{2}}+{\rm j}\biggl(\frac{{\rm i}V^\dagger-U}{\sqrt{2}}\biggr)
		\end{cases}
	\end{equation*}
	for $(Q,P)$, $(X,Y)$, and $(U,V)$ pairs of matrices in $\Hom(\C^m,\C^n)\times\Hom(\C^n,\C^m)$, we establish isometric holomorphic symplectomorphisms between $T^*_{1,0}\Hom(\C^m,\C^n)$ and $\Hom(\Q^m,\Q^n)$ for the structures $(I,\Omega_1)$, $(J,\Omega_2)$, and $(K,\Omega_3)$, respectively. Furthermore, these intertwine the ${\rm U}(m)\subset {\rm Sp}(m)$-action with the ${\rm U}(m)\subset {\rm GL}_m\C$-action on $T^*_{1,0}\Hom(\C^m,\C^n)$.
\end{Lemma}
\begin{proof}
	It is straightforward to verify that the metric in \eqref{metric} is pulled back to \eqref{quat_metric}. One can then show that for the three choices of variables above, the roles of $I$, $J$, and $K$ are cyclically permuted, and thus, so are the holomorphic symplectic forms $\Omega_1$, $\Omega_2$, and $\Omega_3$.
\end{proof}

From now on, suppose $m\le n$ and let $M$ denote the open subset of $\Hom(\Q^m,\Q^n)$ consisting of all elements $A=Q+{\rm j}P^\dagger$ for which $Q$ has maximal rank. The ${\rm U}(m)$-action on $M$ admits a~triple of momentum maps $\mu_k\colon M\rightarrow\mathfrak{u}(m)^*$ for each K\"{a}hler form $\omega_k$. This allows us to introduce the hyperk\"{a}hler quotient
\[
\tilde{M}=\bigl(\mu_1^{-1}({\rm i}\cdot\Id_m)\cap\mu_2^{-1}(0)\cap\mu_3^{-1}(0)\bigr)/{\rm U}(m).
\]
\begin{Theorem}\label{HK_red_thm}
	Let $\bigl(\tilde{M},g,I,J,K\bigr)$ be the hyperk\"{a}hler reduced space for the action of ${\rm U}(m)\subset {\rm Sp}(m)$ on $M\subset\Hom(\Q^m,\Q^n)$ at the regular value $({\rm i}\cdot\Id_m,0,0)\in\mathfrak{u}(m)^*\otimes\R^3$. We have the~following holomorphic symplectomorphisms:
	\begin{itemize}\itemsep=0pt
		\item $\bigl(\tilde{M}, I, \Omega_1\bigr)\cong (T^*_{1,0}\Gr_\C, {\rm i}, \Omega_{\rm can})$,
		\item $\bigl(\tilde{M},J,\Omega_2\bigr)\cong({\rm Orb}(-1), {\rm i}, \Omega_{\rm KKS})$,
		\item $\bigl(\tilde{M},K,\Omega_3\bigr)\cong({\rm Orb}({\rm i}), {\rm i}, \Omega_{\rm KKS})$.
	\end{itemize}
	Here ${\rm Orb}(\zeta)$ is the coadjoint orbit in $\mathfrak{gl}_n\C^*$ equipped with the Kostant--Kirilov--Souriau form $\Omega_{\rm KKS}$ and $\Gr_\C$ is the Grassmannian of complex $m$-dimensional subspaces in $\C^n$.
\end{Theorem}

\begin{proof}
	For the distinguished complex structure $I$, we use Lemma~\ref{facilitate} to identify $\Hom(\Q^m,\Q^n)$ with ${T^*_{1,0}\Hom(\C^m,\C^n)}$. The ${\rm U}(m)$-action admits a holomorphic extension to the ${\rm GL}_m\C$-action demonstrated in Section~\ref{theexample} with momentum map $\mu_m=\mu_2+{\rm i}\mu_3$ given by $\mu_m(Q,P)=PQ$.
	
	In this basis $\mu_1(Q,P)={\rm i}\bigl(Q^\dagger Q-PP^\dagger\bigr)$. We claim that each ${\rm GL}_m\C$-orbit in $M$ intersects $\mu_1^{-1}({\rm i}\cdot\Id_m)$ precisely in a single orbit of ${\rm U}(m)$. To show this, consider how the action of $g$ in ${\rm GL}_m\C$ sends $Q^\dagger Q-PP^\dagger$ to
	\begin{equation}\label{svd}
		g^{-\dagger}Q^\dagger Qg^{-1}-gPP^\dagger g^{\dagger}.
	\end{equation}
	If $Q$ has maximal rank, we may act by the appropriate $g$ to assume that $Q^\dagger Q$ is the identity. If~we~then use the singular value decomposition $g=uDv^\dagger$ and suppose that $v^\dagger PP^\dagger v$ is a~diagonal matrix $\Lambda$ with non-negative diagonal entries, then \eqref{svd} becomes $u\bigl(D^{-\dagger}D^{-1}-D\Lambda D^\dagger\bigr)u^\dagger$. We~can always find a $D$ for which this equals the identity, and hence, have established that the intersection is non-empty.
	
	To show that this intersection is a single ${\rm U}(m)$-orbit, let $\mu_1(Q,P)={\rm i}\cdot\Id_m$ and suppose that~\eqref{svd} is equal to $\Id_m$. By again writing $g$ as $uDv^\dagger$, we find from substituting ${Q^\dagger Q\!=\Id_m\!+\!PP^\dagger}$ into \eqref{svd} that we must have $D^{-\dagger}(\Id_m+S)D^{-1}=\Id_m+DSD^\dagger$, where $S=v^\dagger PP^\dagger v$. This can only hold when $D$ is the identity, and hence, $g$ is unitary.
	
	From this claim, it follows that $\bigl(\tilde{M},I,\Omega_1\bigr)$ may be identified with the holomorphic symplectic reduced space
	\begin{equation*}%\label{QP_quotient}
		\mu_m^{-1}(0)/{\rm GL}_m\C.
	\end{equation*}
	There is a well-defined identification between this quotient and $T^*_{1,0}\Gr_\C$ which sends the equivalence class $[(Q,P)]$ to the plane $\Pi=\Imag Q$ and covector $QP$, interpreted as a map $\C^n/\Pi\rightarrow\C^n$. Observe that the canonical one-form on the cotangent bundle pulls back under this identification to the canonical one-form on $T^*_{1,0}\Hom(\C^m,\C^n)$.
	
	For a different choice of distinguished complex structure, we apply the same argument. For~$J$, we must consider $\mu_m^{-1}(-\Id_m)$ for $\mu_m=\mu_3+{\rm i}\mu_1$, and for $K$ we have $\mu_m^{-1}({\rm i}\cdot\Id_m)$ where ${\mu_m=\mu_1+{\rm i}\mu_2}$. In both cases, the momentum conditions imply that the matrices $(X,Y)$ and~$(U,V)$ all have maximal rank, and a slight alteration to the claim above establishes that the quotients may be identified with the holomorphic symplectic reduced spaces in \eqref{coad_orbit}.
\end{proof}
